% results.tex - Neural Networks format

\section{Results}
\label{sec:results}

\subsection{Main Results on CASME II}

Table~\ref{tab:main_results} presents Censor's performance on CASME II 4-class LOSO evaluation.

\begin{table}[htbp]
\centering
\caption{CASME II 4-class LOSO Results (happiness, surprise, disgust, repression)}
\label{tab:main_results}
\begin{tabular}{lcc}
\toprule
Method & Accuracy & F1-Score \\
\midrule
LBP-TOP~\cite{zhao2007lbptop} & 0.65 & --- \\
Off-ApexNet~\cite{wang2015offapexnet} & 0.68 & --- \\
STSTNet~\cite{ststnet2021} & 0.74 & --- \\
GRA-NET~\cite{graNet2021} & 0.80 & 0.78 \\
MERM~\cite{merm2023} & --- & 0.84 \\
DRConv~\cite{drconv2024} & --- & 0.85 \\
\textbf{Censor (Ours)} & \textbf{0.85} $\pm$ 0.17 & \textbf{0.80} $\pm$ 0.22 \\
\bottomrule
\end{tabular}
\end{table}

Censor achieves 85\% accuracy and 0.80 F1-score, demonstrating competitive performance. The standard deviation ($\pm$0.17) reflects subject variability inherent to LOSO evaluation. Per-fold results range from 0.33 (difficult subject) to 1.00 (ideal subject), with most folds achieving 0.70--0.90.

\subsection{Cross-Dataset Generalization}

Table~\ref{tab:cross_dataset} reports zero-shot transfer results using 3 shared classes.

\begin{table}[htbp]
\centering
\caption{Cross-Dataset Generalization (3-class: happiness, surprise, disgust)}
\label{tab:cross_dataset}
\begin{tabular}{lccc}
\toprule
Source $\rightarrow$ Target & CASME II & SMIC & SAMM \\
\midrule
CASME II $\rightarrow$ & 0.84 & 0.74 & 0.75 \\
SMIC $\rightarrow$ & 0.76 & 1.00 & 0.70 \\
SAMM $\rightarrow$ & 0.67 & 0.69 & 1.00 \\
\bottomrule
\end{tabular}
\end{table}

CASME II $\rightarrow$ SMIC/SAMM achieves 0.74--0.75 accuracy, demonstrating strong cross-dataset generalization. Transfer from larger datasets (CASME II) to smaller ones (SMIC, SAMM) outperforms reverse direction, consistent with expectation that larger training sets yield more generalizable features.

\subsection{Ablation Study}

Table~\ref{tab:ablation} presents component contribution analysis.

\begin{table}[htbp]
\centering
\caption{Ablation Study on CASME II 4-class LOSO}
\label{tab:ablation}
\begin{tabular}{lcc}
\toprule
Configuration & Accuracy & F1-Score \\
\midrule
Fast-only (3D ResNet-18 on optical flow) & 0.86 $\pm$ 0.20 & 0.84 $\pm$ 0.24 \\
Slow-only (3D Swin-T on RGB+rPPG) & 0.67 $\pm$ 0.30 & 0.59 $\pm$ 0.31 \\
Dual + NoMoE & [not available] & [not available] \\
\textbf{Full (Censor: Dual + MoE)} & \textbf{0.85} $\pm$ 0.17 & \textbf{0.80} $\pm$ 0.22 \\
\bottomrule
\end{tabular}
\end{table}

Ablation results reveal an unexpected finding: the fast pathway alone achieves comparable performance to the full model (0.86 vs 0.85 accuracy). This indicates that optical flow features are the primary contributor to micro-expression recognition, capturing the subtle facial muscle movements that define MEs. The slow pathway alone performs significantly worse (0.67 accuracy), suggesting RGB+rPPG features require complementary motion information. The dual-pathway fusion achieves similar performance to fast-only, indicating potential for improved fusion strategies in future work.

\subsection{Analysis}

\subsubsection{Per-Fold Performance Distribution}

LOSO results exhibit substantial variance due to subject heterogeneity. Subjects with fewer samples or unusual facial morphologies yield lower fold accuracy. This variance is inherent to spontaneous micro-expression datasets and motivates LOSO as the standard evaluation protocol.

\subsubsection{Cross-Dataset Transfer Patterns}

SAMM $\rightarrow$ other datasets shows lower performance (0.67--0.69), attributed to SAMM's lower micro-expression intensity and different acquisition settings. CASME II's higher temporal resolution (200 fps) and larger sample count yield more robust training representations.